\documentclass[report.tex]{subfiles}

\begin{document}

\section{Relational Separation Logic}

Having defined our free simulation framework
and instantiated it with \muRTL{},
we require a program logic to specify and verify relational properties
of program executions.
Reasoning about memory transformations using global heap predicates
imposes a heavy verification burden,
as every intermediate step must maintain global invariants.
To enable local reasoning over pairs of memory heaps,
we develop a relational separation logic.

\subsection{Separation Logic}

Separation logic~\parencite{chargueraud-20-sl} establishes
a foundation for modular reasoning about programs that manipulate memory.
In standard Hoare logic,
assertions describe properties of the entire heap.
Establishing that an operation modifies only a restricted portion of memory
requires explicitly proving that all untouched heap regions remain invariant,
an issue known as the frame problem.
Separation logic solves this challenge by introducing
the separating conjunction $P \sep Q$,
which asserts that the memory can be split into two disjoint sub-heaps
satisfying $P$ and $Q$,
respectively.

In verified compilation,
we reason relationally about two program executions:
the target program and the source program.
Because each execution operates over its own memory heap,
we formulate a relational separation logic~\parencite{yang-07-rsl}
whose assertions describe pairs of memory states $(m_t, m_s)$.

In modern interactive proof assistants,
the Iris framework~\parencite{jung-18-iris,krebbers-17-essence} serves
as the standard platform for separation logic in Rocq.
Iris provides expressive mechanisms for higher-order concurrent reasoning,
including step-indexed models,
higher-order ghost state,
and user-defined resource algebras.
For our setting,
these advanced features are not needed.
Instead,
we formalize our relational separation logic directly
as a shallow embedding in Rocq.
We define the type of relational propositions,
denoted $\rProp$,
as a binary relation over memories:
\begin{equation*}
  \rProp \coloneqq \mem \to \mem \to \Prop
\end{equation*}
By convention,
the first argument denotes the target memory $m_t$,
and the second argument denotes the source memory $m_s$.
An assertion $P \in \rProp$ represents a predicate that
expresses ownership over the memory footprints of both programs,
specifying the contents of sub-heaps of the target and source.

\begin{figure}[htb]
  \centering
  \begin{tabular}{lll}
    \toprule
    Connective
    & Semantics
    & Description \\
    \midrule
    $\purel{\psi}$
    & $\lambda m_t\, m_s.\; m_t = \emptyset \land m_s = \emptyset \land \psi$
    & Pure assertion \\
    $\emp$
    & $\lambda m_t\, m_s.\; m_t = \emptyset \land m_s = \emptyset$
    & Empty memory heaps \\
    $P \land Q$
    & $\lambda m_t\, m_s.\; P(m_t, m_s) \land Q(m_t, m_s)$
    & Conjunction \\
    $P \lor Q$
    & $\lambda m_t\, m_s.\; P(m_t, m_s) \lor Q(m_t, m_s)$
    & Disjunction \\
    $P \to Q$
    & $\lambda m_t\, m_s.\; P(m_t, m_s) \implies Q(m_t, m_s)$
    & Implication \\
    $\forall x,\, P(x)$
    & $\lambda m_t\, m_s.\; \forall x,\, P(x)(m_t, m_s)$
    & Universal quantification \\
    $\exists x,\, P(x)$
    & $\lambda m_t\, m_s.\; \exists x,\, P(x)(m_t, m_s)$
    & Existential quantification \\
    $P \sep Q$
    & $\begin{aligned}[t]
      \lambda m_t\, m_s.\; \exists m_t^P m_t^Q m_s^P m_s^Q,\;
      & m_t = m_t^P \uplus m_t^Q \land m_s = m_s^P \uplus m_s^Q  \\
      & {} \land P(m_t^P, m_s^P) \land Q(m_t^Q, m_s^Q)
      \end{aligned}$
    & Separating conjunction \\
    $P \wand Q$
    & $\begin{aligned}[t]
        \lambda m_t\, m_s.\; \forall m_t' m_s',\;
        & P(m_t', m_s') \implies Q(m_t \uplus m_t', m_s \uplus m_s')
      \end{aligned}$
    & Magic wand \\
    $\boxm P$
    & $\lambda m_t\, m_s.\; P(\emptyset, \emptyset)$
    & Persistent modality \\
    \midrule
    $P \ent Q$
    & $\forall m_t\, m_s,\; P(m_t, m_s) \implies Q(m_t, m_s)$
    & Logical entailment \\
    $P \dashv\vdash Q$
    & $(P \ent Q) \land (Q \ent P)$
    & Logical equivalence \\
    \bottomrule
  \end{tabular}
  \caption{Base connectives of our relational separation logic.}
  \label{fig:rprop-connectives}
\end{figure}

The semantic definitions of our base logical connectives are summarized in
\cref{fig:rprop-connectives}.
The entailment $P \ent Q$ asserts that whenever $P$ holds on a memory pair,
$Q$ holds as well.
Logical equivalence $P \dashv\vdash Q$ holds
when both directions of entailment hold.
The connectives satisfy the standard algebraic properties of separation logic:

\begin{itemize}
  \item The separating conjunction $P \sep Q$ splits both the target heap $m_t$
        and the source heap $m_s$ into disjoint sub-heaps%
        \footnote{Here $\uplus$ denotes the disjoint union of memory heaps,
          which is defined only when its operand domains are disjoint.},
        assigning one part to $P$
        and the other part to $Q$.
        Because the sub-heaps must be strictly disjoint,
        this operator prevents duplicate resource ownership.

  \item The persistent modality $\boxm P$ asserts that $P$ holds on empty heaps.
        Because a persistent assertion requires no memory resources,
        it represents a duplicable fact:
        $\boxm P \ent \boxm P \sep \boxm P$.
\end{itemize}

Our memory logic is \emph{linear} rather than affine.
In an affine logic,
heap chunks can be silently discarded via weakening:
$P \sep Q \ent P$.
In our linear formulation,
proving $P \ent Q$ guarantees that every memory cell described by $P$
is accounted for in $Q$.
This strict resource discipline forbids any memory leak.

To preserve the interactive proof engineering benefits of the
Iris Proof Mode (MoSeL)~\parencite{krebbers-18-mosel},
we prove that our connectives satisfy the required typeclass hierarchy.
This allows our formalization to use standard Iris tactics
while executing directly on our shallow relational model.

\subsection{Relational Assertions \& Memory Predicates}
\label{sec:mem-predicate}

Building on our memory model (\cref{fig:mem-and-reg}),
we introduce primitive memory assertions for target and source states.
We define relational points-to and freed assertions:
\begin{align*}
  \ptrt{l}{v}
  &
    \coloneqq
    \lambda m_t\, m_s.\;
    m_t = \{ l \mapsto \Allocated{v} \} \land m_s = \emptyset
  \\
  \ptrs{l}{v}
  &
    \coloneqq
    \lambda m_t\, m_s.\;
    m_t = \emptyset \land m_s = \{ l \mapsto \Allocated{v} \}
  \\
  \freet{l}
  &
    \coloneqq
    \lambda m_t\, m_s.\;
    m_t = \{ l \mapsto \Freed \} \land m_s = \emptyset
  \\
  \frees{l}
  &
    \coloneqq
    \lambda m_t\, m_s.\;
    m_t = \emptyset \land m_s = \{ l \mapsto \Freed \}
\end{align*}
The predicate $\ptrt{l}{v}$ asserts that location $l$ in the target heap
contains value $v$,
while requiring the source heap to be empty.
Symmetrically,
$\ptrs{l}{v}$ asserts ownership of an allocated cell in the source heap.
The predicates $\freet{l}$ and $\frees{l}$ assert ownership
of deallocated memory cells in the target and source heaps,
respectively.

Memory allocations in the target and source programs
are not guaranteed to produce identical physical addresses.
To account for address variations,
we maintain an address mapping $I \subseteq \loc \times \loc$,
where each pair $(l_t, l_s) \in I$ relates a target location $l_t$
to a source location $l_s$.
We enforce that $I$ is a bijective mapping:
target locations are uniquely related to source locations.
This mapping $I$ parameterizes the value similarity relation
$\sameval{I}{v_t}{v_s}$ (\cref{fig:same-val}).
As the bijection grows monotonically with fresh memory allocations,
the value relation is preserved:
\begin{lemma}[Monotonicity of $\sameval{I}{}{}$]
  \label{lem:sameval-mono}
  If $I \subseteq I'$,
  then $\sameval{I}{v_t}{v_s} \implies \sameval{I'}{v_t}{v_s}$.
\end{lemma}

To relate entire target and source heaps across program execution,
we adapt the memory injection predicate $\meminj{I}{E}$
defined by \textcite{gaher-22-simuliris}.
When the set $E$ is empty,
this predicate owns all memory cells registered in $I$.
It asserts that every active pair of locations $(l_t, l_s) \in I$
holds relationally similar values,
and that deallocated cells are consistently marked as freed in both heaps.

When verifying an instruction that reads,
writes,
or deallocates a specific pair of locations $(l_t, l_s) \in I$,
we need to access the underlying points-to or freed assertion.
However,
removing $(l_t, l_s)$ from the mapping $I$ would shrink $I$,
breaking its monotonicity
and invalidating the value similarity of existing pointers.
To overcome this difficulty,
we introduce the \emph{exploited set} $E \subseteq I$.
The set $E$ tracks the address pairs whose individual heap assertions
have been temporarily extracted from the global injection for local reasoning.

The interaction between the memory injection predicate,
the bijection $I$,
and the exploited set $E$ is primarily governed by three lemmas:
\begin{itemize}
  \item \textbf{Exploitation:}
        When $(l_t, l_s) \in I$ is unexploited,
        the assertion $\meminj{I}{E}$ can be opened to extract ownership
        of $(l_t, l_s)$ while updating the exploited set to $E \cup \{ (l_t, l_s) \}$.

  \item \textbf{Reintegration:}
        Once a memory operation completes,
        the updated points-to assertion $\ptrt{l_t}{v_t} \sep \ptrs{l_s}{v_s}$
        (with $\sameval{I}{v_t}{v_s}$)
        or freed assertion $\freet{l_t} \sep \frees{l_s}$
        can be released back into the injection,
        decreasing the exploited set to $E \setminus \{ (l_t, l_s) \}$.

  \item \textbf{Extension:}
        When allocating fresh memory cells,
        the fresh pair of addresses $(l_t, l_s)$
        can be added to the memory injection
        provided their values are related by $\sameval{I}{}{}$.
\end{itemize}

\subsection{Relational Hoare Quadruples}
\label{sec:rel-hoare}

By combining the free simulation with relational separation logic,
we define relational Hoare quadruples.
In fixed target and source programs $P_t$ and $P_s$,
a relational Hoare quadruple relates a target function $f_t$
to a source function $f_s$.
The formal definition is given by:
\begin{equation*}
  \begin{aligned}
    \hoareq{P}{f_t}{f_s}{Q} \coloneqq
    \boxm \Big(
    \forall \vec{v}_t\, \vec{v}_s,\;
    P(\vec{v}_t, \vec{v}_s) \wand
    \fsim{([], \CallState{f_t}{\vec{v}_t})}{0}{0}{([], \CallState{f_s}{\vec{v}_s})}{Q}
    \Big)
  \end{aligned}
\end{equation*}
where $P$ is a precondition relating the target and source argument vectors,
and $Q$ is a postcondition referring to the returned values.
The simulation is here lifted to an assertion in $\rProp$:
\begin{equation*}
  \fsim{t}{j}{i}{s}{\Phi}
  \coloneqq
  \lambda m_t\, m_s.\;
  \fsim{(t, m_t)}{j}{i}{(s, m_s)}{\Phi}
\end{equation*}


The meaning of this quadruple is given by the free simulation:
whenever the initial function arguments and heap states satisfy
$P(\vec{v}_t, \vec{v}_s)$,
if the source execution terminates without encountering undefined behavior,
then the target execution also terminates with return values and memories
satisfying $Q(v_t, v_s)$.
If the source execution diverges,
the target execution diverges as well.
If the source execution reaches a stuck state,
the simulation relation grants the compiler complete freedom.

The definition of the Hoare quadruple
is guarded by a persistent modality $\boxm$
to guarantee that the specification holds independently of heap ownership.
As a consequence,
a verified quadruple constitutes a duplicable resource that can be reused
across multiple call sites throughout the program.

\end{document}

% LocalWords:  Coinductive coinductive asymmetricity arities
